\section{Python 数值实验}\label{ch:experiments}

前面五课都是先通过手算理解，再逐渐引入通用记号。现在使用 Python 把一些计算重复很多次，并核对结果。所有示例都可以从仓库根目录运行。先执行 \texttt{python -m pip install -r requirements.txt} 安装 NumPy。

\subsection{网格搜索与参数估计}

第一课的三个训练样本是 $(1,2),(2,4),(3,6)$，模型为 $\hat y=wx$。我们知道 $w=2$ 时 MSE 为零。现在不直接用解析公式，而是让计算机枚举 $0$ 到 $4$ 之间的一些候选参数，每次计算
\[
J(w)=\frac13\big[(w-2)^2+(2w-4)^2+(3w-6)^2\big].
\]
代码保存在 \texttt{code/grid\_search\_intro.py}：
\lstinputlisting[caption={从损失函数出发搜索最佳参数},label={lst:grid-search}]{code/grid_search_intro.py}

它会输出最优候选参数 $2.0$、最小训练 MSE 为 $0.0$，以及 $x=4$ 时的预测 $8.0$。网格搜索只在\emph{列举过的参数}中找最优值；本例恰好包含真实最优参数，并不意味着一般情形也会找到精确最优解。

\subsection{最小二乘的 NumPy 实现}

第二课使用四组数据 $(1,2),(2,3),(3,5),(4,4)$。先求 $\bar x,\bar y$，再计算
\[
\hat w=\frac{\sum_i(x_i-\bar x)(y_i-\bar y)}{\sum_i(x_i-\bar x)^2},
\qquad
\hat b=\bar y-\hat w\bar x.
\]
结果应该是 $\hat w=0.8,\hat b=1.5$，训练 MSE 为 $0.45$。这次不必依赖封装好的回归库：把数学公式直接写成 \texttt{np.mean} 与 \texttt{np.sum}，可以看到每一步对应哪条公式。示例代码：
\lstinputlisting[caption={直接将最小二乘公式翻译为 NumPy},label={lst:ols-intro}]{code/ordinary_least_squares_intro.py}

仓库另外保留 \texttt{code/linear\_regression\_demo.py}，用于比较最小二乘数值解与梯度下降结果。原有脚本使用另一组数据 $(1,2),(2,3),(3,5)$：这组数据的斜率是 $3/2$、截距是 $1/3$，不是刚才四样本例子的 $0.8$、$1.5$。\textbf{不要混用不同实验的数据和答案。}

\subsection{梯度下降的程序实现}

第三课已经手算过
\[
f(x)=(x-3)^2,\qquad x_0=1,\qquad\eta=0.2.
\]
第一次更新得到 $x_1=1.8$，之后是 $2.28$、$2.568$ 等。使用一个 \texttt{for} 循环，把“求导 $\rightarrow$ 更新 $\rightarrow$ 计算新函数值”重复即可：
\lstinputlisting[caption={把一元函数梯度下降写成循环},label={lst:gd-intro}]{code/gradient_descent_intro.py}

在更复杂的线性回归里，循环内部只需改为同时更新 $w,b$。关键是先从旧参数处计算两个偏导，再一起更新。软件重复操作很快，但理解它的最好方式仍是手算第一步。

\subsection{逻辑回归的梯度更新}

第四课用一个样本 $x=2,y=1$ 演示：
\[
w_0=b_0=0,\quad\eta=0.1,\quad
p_0=\frac1{1+e^0}=0.5.
\]
解析梯度为 $\partial_wJ=(p-y)x=-1$ 与 $\partial_bJ=p-y=-0.5$。所以一步更新得 $w_1=0.1,b_1=0.05$，新概率为 $\sigma(0.25)\approx0.5622$：
\lstinputlisting[caption={一个样本的逻辑回归梯度更新},label={lst:logit-intro}]{code/logistic_regression_intro.py}

仓库原有的 \texttt{code/logistic\_gradient\_check.py} 则做另一件事：用数值差分近似求导，检查解析梯度的实现是否正确。对某个参数分量 $\beta_j$，中心差分近似为
\begin{equation}\label{eq:finite-difference}
\frac{\partial J}{\partial\beta_j}
\approx\frac{J(\boldsymbol\beta+h\boldsymbol e_j)-J(\boldsymbol\beta-h\boldsymbol e_j)}{2h}.
\end{equation}
这里不要求初学者独立证明该近似，只需要知道它可以帮助发现代码求导错误。

\subsection{多项式拟合与过拟合}

第五课介绍多项式复杂度与过拟合。仓库已有 \texttt{code/polynomial\_comparison.py}，可直接运行，记录不同次数模型的训练、测试 MSE。在比较结果时，一定看清数据生成方式、训练与测试的划分，以及数值是否来自同一随机种子。

\subsection{实验可复现性}

为了真正复现，一份实验记录至少需要交代：数据或生成规则、样本数量、训练/验证/测试划分、随机种子、模型形式、损失函数、优化方法、学习率或其他超参数、评价指标、Python 与依赖版本。

\textbf{这部分的主要目标不是堆积代码，而是建立“公式 $\leftrightarrow$ 手工计算 $\leftrightarrow$ Python 结果”的联系。}如果输出与手算不一致，优先核对数据是否相同、分母是 $n$ 还是 $2n$、残差方向以及更新时有没有混用新旧参数。
